% Commentary and audit by GPT-6 Astra (OpenAI, ChatGPT Pro mode).
\begin{otherlanguage*}{english}
\noindent\textbf{Presentations which retain structure.}
These pages turn the existence of small approximations into a method for handling objects with structure. An accessible category supplies filtered presentations of its objects; a diagram requires compatible presentations of its arrows and their equations. A section of a fibred category adds transport between different fibres. A full subcategory defined by a property requires approximations which already possess that property. The successive arguments address these problems with the same operations: choose a common stage, impose a bounded family of equations, and take a suitably long colimit. The cardinal estimates specify when this process can still be carried out among the chosen small objects.

These pages begin with Theorem 4.9.3 and end during the proof of Theorem 4.13.15, on manuscript page 328 \citep{Grothendieck1990XVIII}. They are printed pages 79--123 of the edition; the manuscript labels run from an opening continuation before 220 through 328, including 302 bis and the cancelled page 320. These systems should not be identified with one another. The complete chapter was consulted where it supplies the opening hypotheses or the continuation of an unfinished proof. The end of these pages is not itself a mathematical gap.

\medskip\noindent\textbf{The cardinal convention.}
Grothendieck's phrase ``grande devant $\pi$'' means modern $\pi^+$-filteredness: every diagram with at most $\pi$ arrows has a cocone. His $\pi$-accessible objects are therefore $\pi^+$-presentable objects, and his $\Ind_\pi C$ is the completion under $\pi^+$-filtered colimits. This translation matters even for $\pi=\aleph_0$, where countable constraints, rather than only finite constraints, are admitted. It also prevents an unnoticed assumption that $\pi$ is regular. The adjective ``regular'' applied to a category's filtration elsewhere in the chapter denotes a separate categorical condition.

A second distinction concerns a full subcategory $\mathcal S\subseteq\mathcal M$. Intrinsic accessibility, accessibility of the inclusion, and preservation of presentable objects by that inclusion are different requirements. The source's fixed-rank notion of accessibility \emph{in} $\mathcal M$ requires the last comparison as well. We keep the intrinsic notation $\mathcal S_\pi$ separate from $\mathcal S\cap\mathcal M_\pi$ until the hypotheses identify them.

\medskip\noindent\textbf{Cofinality and simultaneous presentation.}
The first task is to show that a small diagram in $\Ind_\pi C$ comes from a filtered system of diagrams in $C$. Full faithfulness follows from expressing natural transformations as an equalizer of products of hom sets, provided the indexing category has at most $\pi$ arrows. Essential surjectivity is the substantial part. The source supplies it when $C$ has colimits of one suitable limit-ordinal shape. In the ensuing section theorem, that colimit turns approximately compatible structure maps into actual structure maps.

The cofinality discussion 4.9.6 isolates the mechanism. We use ``final'' for right-cofinal functors. A faithful functor can reflect an equality once the requisite arrows have been lifted after enlargement. Without faithfulness, eventual equalization of arrows having the same image suffices; the commentary proves this relative variant explicitly. It also completes the cardinal descent and the discrete-opfibration lemma. Common upper bounds, connectedness, and equalization are kept distinct, with worked examples showing why none can silently replace another.

One clause is missing from the printed Lemma 4.9.6.3(b$'$): filteredness of the comma categories. The faithful functor from a two-object discrete category to the terminal category disproves the formulation as printed. The proof on page 236 specifies the missing condition. The possible empty-category issue in b$''$ is instead a convention to clarify: its cardinal version already includes the empty family among the families to be bounded.

The finite-diagram question now has a definitive answer. The comparison works for every $C$ exactly for essentially small enough well-founded shapes, in the sense of \citet[Theorem 1.3]{Henry2025Ind}. The finite rigid case in 4.11.6 is included. A nontrivial finite group supplies an obstruction, and the commentary calculates the signed-ordinal involution example at every infinite completion bound. This is a known construction, with its antecedent in SGA~4, I, 8.8.8, not a result claimed for the commentators \citep{SGA4I1972}.

\medskip\noindent\textbf{Sections and the transfer of colimits.}
In 4.10 ``cofibrant'' describes a Grothendieck opfibration. Its sections are families $X_i$ with maps $a_*X_i\to X_j$ for $a:i\to j$; those maps need not be cocartesian. The approximation category of sections over a given section maps finally to the product of the fibrewise approximation categories. Evaluating its canonical colimit therefore recovers the original section. We write out this passage, which the source asks to have made explicit, and prove full faithfulness separately.

The transfinite construction has a coherence defect as displayed. Compatibility of consecutive stages does not compare choices across an internal limit ordinal. A constant two-point diagram whose involution changes at $\omega$ satisfies the displayed conditions but gives no induced map on its constant colimit. The repair imposes compatibility with every earlier stage when constructing a new stage; there are still at most $\pi$ equations. The completed induction proves the intended section theorem. A separate endpoint adjustment handles the parenthesis allowing an auxiliary cardinal equal to $\pi$.

The application to paratoposes concerns locally presentable categories, without an additional assertion of topos exactness. Ordinary sections of an opfibration have pointwise limits whenever the fibres have them: their structure maps are induced by the limit projections. Preservation by transport is needed for pointwise colimits. This distinction vindicates the printed limit argument in 4.10.9 and prevents the editors' tentative alternative from becoming an unwarranted correction.

Question 4.10.11 about ordinary and cartesian sections is settled by the later theory of accessible limits. Products, inserters, and equifiers encode arbitrary sections; requiring invertible structure maps gives the pseudolimit describing cartesian sections. After choosing common ranks, \citet[Theorems 8.1, 8.3 and Corollary 9.2]{Positselski2024Limits} apply. Their fixed-rank conclusions retain hypotheses about a smaller chain-colimit cardinal, close to the distinction Grothendieck is already exploiting.

Section 4.11 transfers bounded colimits from $C$ to $\Ind_\pi C$ by first presenting the entire diagram. The universal-property argument interchanges two limits of hom sets. The separate assertion that $C\to\Ind_\pi C$ preserves existing bounded colimits uses filtered-colimit/limit commutation in sets. Idempotent completeness enters in the converse, through the description of presentable ind-objects as retracts of objects of $C$. We also give the finite rigid induction and the reduction of finite limits and colimits to products and equalizers, or coproducts and coequalizers.

\medskip\noindent\textbf{Accessible functors and accessible properties.}
Uniformization in 4.12 bounds the ranks of the images of a small generating family. At a sufficiently large fixed point $\pi^{\pi_0}=\pi$, an accessible functor carries $\pi$-accessible objects to $\pi$-accessible objects. Fixing a preservation rank reduces the functor category to one with small domain. The proof of 4.12.4 invokes the ind universal property without explicitly securing the required colimits in the target. Choosing the rank useful for both categories repairs the proof, as in 4.12.5. We distinguish this verified form from the broader first assertion rather than manufacture a counterexample from an omitted hypothesis.

There is no common rank for all accessible endofunctors of sets: the representables $X\mapsto X^B$ prove this as $|B|$ grows. Small colimits of accessible set-valued functors are pointwise, so the same example rules out a small accessible generating family for their entire category. This supplies the counterexample sought in Remark 4.12.6.

The existential theorems about accessible full subcategories and their small intersections remain valid. A later higher-categorical counterpart is \citet[Proposition 5.4.7.10]{Lurie2009HTT}. The effective threshold in 4.13.2(c), however, omits the ambient ranks of the generators. The singleton subcategory $\{X\}\subseteq\mathcal P(X)$ gives a complete counterexample to that estimate. The repair retains the missing bound already used on page 282.

For intersections, an ordinal of cardinality $\pi$ need not have cofinality greater than the smaller bound. We prove a replacement construction which attends to all the subcategories at each stage of a chain of regular length $\rho^+\leq\pi$. Its colimit belongs to every subcategory. This proves the main cofinal-intersection lemma without an unnecessary regularity restriction on $\pi$. A related argument repairs the cofinal-pullback construction in the final proof.

Theorem 4.13.15 needs a more substantial correction. Under its printed hypothesis $\pi^{\pi_0}=\pi$, condition a$'$ need not imply a): closure of the presentable part under $\pi_0'^+$-filtered colimits of size at most $\pi$ need not give unrestricted closure. We construct a counterexample with $\pi_0=\aleph_0$, $\pi_0'=\aleph_1$, and $\pi=\beth_{\omega_1}$, proving both the bounded closure and the failure of unrestricted closure. Strengthening the hypothesis to $\pi^{\pi_0'}=\pi$ gives a sufficient repair: the indexing categories can then be exhausted by suitably filtered subcategories of size at most $\pi$. The bounded membership criteria and the accessible-intersection theorem survive this failed implication.

\medskip\noindent\textbf{Editorial and mathematical status.}
The registers distinguish copying errors from source slips, and false inferences from conclusions with repaired proofs. Confirmed source errors remain visible in the text. The results presented as our developments are marked \ours; completed source arguments and known constructions are distinguished explicitly. The literature checked is stated without a priority claim. This commentary and audit are AI-produced by GPT-6 Astra (OpenAI, ChatGPT Pro mode), checked and integrated by Claude (Anthropic).
\end{otherlanguage*}
